\scanpage{212}
cette condition sous la forme (29'),

\[
\unsure{(*)}\quad\forall s,t\in S,\quad\{s,t\}\in A_0\iff\{\lambda^{-1}s,\lambda t\}\in A_0,
\]

on voit ceci

\[
\begin{aligned}
&\text{Si }\lambda,\mu\in\unsure{\mathfrak S|\mathscr S}_S\ \text{satisfont séparément à la condition}\\
&\unsure{(*)|\text{cas lin.}}\ \text{(i.e. sont des homothéties pour le graphe }(S,A_0)),\\
&\text{alors }\lambda\mu\ \text{y satisfait \unsure{sss|ssi|si} on a }\lambda\mu=\mu\lambda\\
&\deleted{\text{\unsure{sur}}}\quad\text{(cas d'un graphe \uline{séparant}).}
\end{aligned} \tag*{(66)}
\]

Pour continuer le travail, on supposera donc donné un sous-groupe \uline{commutatif}

\[
\deleted{\unsure{Z}}\qquad Z\subset\mathfrak S_S
\]

formé d'homothéties de $(S,A_0)$, et on considère \inserted{alors} le \deleted{éléments} \inserted{sous-groupe} correspondant \inserted{$Z_\Gamma$} de \deleted{$\unsure{\mathfrak S_S}\ \mathfrak S_S\times\mathfrak S_{S^\vee}$} $\operatorname{Aut}(S,S^\vee,R)$, on trouve : \deleted{que les}

\[
\left\{\begin{aligned}
&\text{les }\varphi_{\unsure{a|\alpha}}:S\longrightarrow S^\vee\ \text{telles qu'on ait}\\
&\lambda^\vee\varphi_a\lambda=\varphi_a\quad\text{pour tout }\lambda\in Z_\Gamma\subset\operatorname{Aut}_S(\Gamma)\\
&\text{forment un }\deleted{\text{torseur sous le groupe }Z_\Gamma,}\ \inserted{\unsure{Z_\Gamma|\Gamma}\text{-ens }\Pi_Z,}\ \text{par}\\
&\text{par l'opération }\deleted{\text{\unsure{qui à}}}\ \text{du groupe }Z_\Gamma,\ \text{qui}\\
&\hspace{48mm}\text{y opère \uline{librement} par}\\
&T_{\lambda,\lambda^\vee}\varphi=\lambda^\vee\varphi=\varphi\lambda^{-1}.
\end{aligned}\right. \tag*{(67)}
\]

% slanted
\margin{On a supposé que $\Pi_Z\ne\varnothing$, car on est parti d'une $\varphi_0\in\Pi_Z$.}

En effet, si \inserted{$(\lambda,\lambda^\vee)\in Z$,} on doit vérifier \inserted{d'abord} que

\[
\inserted{\varphi\in\Pi_Z}\Longrightarrow\underbrace{T_{\lambda,\lambda^\vee}\varphi}_{\lambda^\vee\varphi}\in Z\quad\text{i.e.}\quad
\forall(\mu,\mu^\vee)\in Z,\ \deleted{\text{qu'}}\text{on a}
\]

\[
\begin{aligned}
&\underbrace{\mu^\vee(\lambda^\vee}_{\lambda^\vee\mu^\vee}\varphi)\mu=(\lambda^\vee\varphi)\\
&\hspace{18mm}\underbrace{\phantom{\mu^\vee\varphi\mu}}_{\varphi}
\end{aligned}
\]

O,K, Le fait que $Z_\Gamma$ opère librement est clair, \deleted{\unsure{pour}} car pour $\varphi_0,\varphi$ \unsure{fixés|fixées}, il n'y a qu'un seul $\lambda,\lambda^\vee$ \ldots

\deleted{On \unsure{fait} le fait} \deleted{ensuite que si $\varphi_0\in$} On aimerait que

% archive sheet number: 211
